
\documentclass{article}
\usepackage{colortbl}
\usepackage{makecell}
\usepackage{multirow}
\usepackage{supertabular}

\begin{document}

\newcounter{utterance}

\centering \large Interaction Transcript for game `eqbench', experiment `eq\_bench', episode 48 with reflexion{-}llama{-}3.1{-}8b{-}instruct.
\vspace{24pt}

{ \footnotesize  \setcounter{utterance}{1}
\setlength{\tabcolsep}{0pt}
\begin{supertabular}{c@{$\;$}|p{.15\linewidth}@{}p{.15\linewidth}p{.15\linewidth}p{.15\linewidth}p{.15\linewidth}p{.15\linewidth}}
    \# & \multicolumn{2}{c}{Player} && \multicolumn{2}{c}{Game Master} \\
    \hline

    \theutterance \stepcounter{utterance}  
    & & & \multicolumn{4}{p{0.6\linewidth}}{
        \cellcolor[rgb]{0.9,0.9,0.9}{
            \makecell[{{p{\linewidth}}}]{
                \texttt{\tiny{[P1$\langle$GM]}}
                \texttt{You have attempted this same task before and did not succeed. These are the plans you wrote after those attempts.} \\
\texttt{Trial \#1: New plan: I failed to accurately predict Dr. Hall's emotional response because I didn't consider the context of the dialogue and the tone of the conversation. Dr. Hall's comment about Jen's wit being "as sharp as a butter knife" was a veiled insult, and Jen's response was a clever comeback. However, I didn't account for the fact that Dr. Hall's comment was a subtle jab, and Jen's response was a more overt attack. Next time, I will pay closer attention to the tone and context of the dialogue, and consider the possibility that Dr. Hall's comment was a backhanded compliment. I will also look for more subtle cues in the dialogue, such as body language or tone of voice, to inform my prediction of Dr. Hall's emotional response.} \\
\\ 
\\ 
\\ 
\texttt{Your task is to predict the likely emotional responses of a character in this dialogue:} \\
\\ 
\texttt{Dr. Hall: "Ms. Vasquez, I see you're gracing us with your presence today. I trust you've done the reading?"} \\
\texttt{Jen: "Oh, I wouldn't miss it for the world, Dr. Hall. And yes, I've done the reading. Twice."} \\
\texttt{Dr. Hall: "Twice? My, my. Overachieving or overcompensating, I wonder?"} \\
\texttt{Jen: "Funny you should mention overcompensation, Dr. Hall."} \\
\texttt{Dr. Hall: "Careful, Ms. Vasquez. Your wit is as sharp as a butter knife."} \\
\texttt{Jen: "And your condescension is as subtle as a sledgehammer, Dr. Hall."} \\
\texttt{{[}End dialogue{]}} \\
\\ 
\texttt{At the end of this dialogue, Dr. Hall would feel...} \\
\texttt{Apologetic} \\
\texttt{Angry} \\
\texttt{Surprised} \\
\texttt{Superior} \\
\\ 
\texttt{Give each of these possible emotions a score from 0{-}10 for the relative intensity that they are likely to be feeling each.} \\
\\ 
\texttt{You must output in the following format, including headings (of course, you should give your own scores), with no additional commentary:} \\
\\ 
\texttt{Apologetic: <score>} \\
\texttt{Angry: <score>} \\
\texttt{Surprised: <score>} \\
\texttt{Superior: <score>} \\
\\ 
\\ 
\texttt{{[}End of answer{]}} \\
\\ 
\texttt{Remember: zero is a valid score, meaning they are likely not feeling that emotion. You must score at least one emotion > 0.} \\
\\ 
\texttt{Your answer:} \\
            }
        }
    }
     \\ \\

    \theutterance \stepcounter{utterance}  
    & \multicolumn{4}{p{0.6\linewidth}}{
        \cellcolor[rgb]{0.561,0.737,0.561}{
            \makecell[{{p{\linewidth}}}]{
                \texttt{\tiny{[P1$\rangle$GM]}}
                \texttt{Apologetic: 0} \\
\texttt{Angry: 8} \\
\texttt{Surprised: 2} \\
\texttt{Superior: 0} \\
            }
        }
    }
    & & \\ \\

    \theutterance \stepcounter{utterance}  
    & & & \multicolumn{2}{p{0.3\linewidth}}{
        \cellcolor[rgb]{0.9,0.9,0.9}{
            \makecell[{{p{\linewidth}}}]{
                \texttt{\tiny{[GM$|$GM]}}
                \texttt{Apologetic: 0} \\
\texttt{Angry: 8} \\
\texttt{Surprised: 2} \\
\texttt{Superior: 0} \\
            }
        }
    }
    & & \\ \\

    \theutterance \stepcounter{utterance}  
    & & & \multicolumn{2}{p{0.3\linewidth}}{
        \cellcolor[rgb]{0.9,0.9,0.9}{
            \makecell[{{p{\linewidth}}}]{
                \texttt{\tiny{[GM$|$GM]}}
                \texttt{Apologetic: 0} \\
\texttt{Angry: 6} \\
\texttt{Surprised: 4} \\
\texttt{Superior: 8} \\
            }
        }
    }
    & & \\ \\

    \theutterance \stepcounter{utterance}  
    & & & \multicolumn{2}{p{0.3\linewidth}}{
        \cellcolor[rgb]{0.9,0.9,0.9}{
            \makecell[{{p{\linewidth}}}]{
                \texttt{\tiny{[GM$|$GM]}}
                \texttt{game\_result = LOSE} \\
            }
        }
    }
    & & \\ \\

    \theutterance \stepcounter{utterance}  
    & & \multicolumn{4}{p{0.6\linewidth}}{
        \cellcolor[rgb]{0.95,0.95,0.95}{
            \makecell[{{p{\linewidth}}}]{
                \texttt{\tiny{[P1$\leftrightarrow$P]}}
                \texttt{Post{-}episode Reflexion (not sent to GM):} \\
\texttt{New plan: I failed to accurately predict Dr. Hall's emotional response because I didn't consider the context of the dialogue and the tone of the conversation. Dr. Hall's comment about Jen's wit being "as sharp as a butter knife" was a veiled insult, and Jen's response was a clever comeback. However, I didn't account for the fact that Dr. Hall's comment was a subtle jab, and Jen's response was a more overt attack. Next time, I will pay closer attention to the tone and context of the dialogue, and consider the possibility that Dr. Hall's comment was a backhanded compliment. I will also look for more subtle cues in the dialogue, such as body language or tone of voice, to inform my prediction of Dr. Hall's emotional response.} \\
\\ 
\texttt{To improve my prediction, I will consider the following:} \\
\\ 
\texttt{1. Identify the tone of the conversation: Is it formal or informal? Is it sarcastic or genuine?} \\
\texttt{2. Analyze the language used: Are there any veiled insults or backhanded compliments?} \\
\texttt{3. Look for subtle cues: Are there any nonverbal cues, such as body language or tone of voice, that can inform my prediction?} \\
\texttt{4. Consider the context: What is the purpose of the conversation? Is it a formal meeting or an informal discussion?} \\
\\ 
\texttt{By considering these factors, I will be able to make a more accurate prediction of Dr. Hall's emotional response.} \\
\\ 
\texttt{Apologetic: 0} \\
\texttt{Angry: 9} \\
\texttt{Surprised: 0} \\
\texttt{Superior: 1} \\
            }
        }
    }
    & \\ \\

\end{supertabular}
}

\end{document}
